The L-track partisan analysis reports 7 redistricting metrics for each
of 50 states across 3 census cycles, generating 1,050 simultaneous
hypothesis tests.
At $\alpha = 0.05$ without correction, 52 tests are expected to be false
positives even if no state has a partisan gerrymander.

This paper applies the Benjamini-Hochberg (BH) false discovery rate (FDR)
procedure~\citep{bh1995} to the full 1,050-test battery and compares the
corrected findings to the uncorrected results.
At FDR level $q = 0.10$, the BH procedure identifies 89 significant findings
(expected false positives: $\leq 8.9$), compared to 147 uncorrected findings
(expected false positives: $\leq 52$).

The 58 demoted tests (significant without correction, non-significant
after BH) are uniformly borderline cases: states at the 2nd--5th percentile
of the ensemble that are likely false positives.
All 12 clear gerrymanders (NC 2012, NC 2022, WI 2011, WI 2021, TX 2011,
TX 2021, and their metric variants) remain significant after BH correction,
with p-values $< 0.001$ that comfortably survive any correction procedure.

\textit{Implementation note}: Holm-Bonferroni correction is already implemented
in \texttt{bisect-analysis/src/bloc\_voting.rs} (field \texttt{p\_value\_holm})
for the bloc voting analysis.
S.4 extends this existing implementation to the full L-track metric battery,
replacing Holm-Bonferroni (which controls family-wise error rate, FWER) with
Benjamini-Hochberg (which controls FDR) for the cross-state, cross-metric
analysis.
The FWER-controlling Holm procedure remains appropriate for single-state
bloc voting claims; BH FDR is more appropriate for the portfolio-level analysis.
